\section{技术专家一：VLA 是物理世界的 Driver OS}

上一章讲 AI 产品和组织方法，本章进入技术专家通道。张小珺调用“CEO 大模型的技术专家”，让李想讲 VLA。VLA 即 Vision-Language-Action：Vision 负责感知物理世界，Language 负责理解规则、意图和知识，Action 负责输出动作。对理想而言，VLA 不是抽象多模态玩具，而是“司机大模型”，也是汽车和交通领域更重要的大模型或操作系统。

\begin{figure}[H]
\centering
\includegraphics[width=0.86\textwidth]{vla-driver-os-stack.png}
\caption{VLA 司机操作系统：Vision、Language、Action 组合成物理世界的驾驶大模型。自制概念图，依据 00:36:18--00:42:30 对谈内容整理。}
\end{figure}

\begin{knowledgebox}{读图：VLA 的三层不是并列名词}
Vision 解决“看见”，Language 解决“理解”，Action 解决“做”。自动驾驶真正难的是把三者放到一个可训练、可评测、可部署的闭环里，而不是分别把视觉模型、语言模型和控制器拼在一起。
\end{knowledgebox}

\subsection{从规则算法到端到端，再到 VLA}

李想把智能驾驶的演化描述为进化过程，而不是突变。早期是规则算法，之后是端到端加 VLM，现在进入 VLA 阶段。这里的关键变化是，系统不再只是“看懂”和“规划”，而是要把动作输出纳入学习。VLA 的目标不是让车学会讲道理，而是让车在真实道路上做出更像人、更稳定、更安全的驾驶动作。

\begin{importantbox}{VLA 的技术定义}
VLA 是把视觉输入、语言/知识表示和动作输出放入统一训练与推理链路的模型。对自动驾驶而言，它的输出不是文本，而是车辆行为；因此它的评测也不能只看回答质量，而要看安全、舒适、规则遵守和目的达成。
\end{importantbox}

\subsection{32B 云端机座与车端部署}

访谈提到，理想正在训练 32B 的云端 VL 机座模型。云端模型负责吃更大数据、更大参数和更复杂训练；车端模型则受限于算力、延迟和安全要求，需要通过蒸馏、部署和验证进入车辆。这个结构和前面 EP120 小鹏谈到的云端模型工厂相互呼应：物理世界 AI 往往不是“一个模型上车”，而是“云端训练、端侧压缩、真实验证”的系统工程。

\begin{knowledgebox}{术语消化：VLA 训练链路}
\begin{tabularx}{\textwidth}{p{0.22\textwidth}p{0.32\textwidth}X}
\toprule
环节 & 作用 & 风险与约束 \\
\midrule
VL 机座 & 让模型理解视觉和语言 & 需要大规模多模态数据和算力。 \\
Action 训练 & 把理解转为控制动作 & 输出直接影响安全和体验。 \\
世界模型数据 & 提供可控训练和测试场景 & 仿真与真实世界之间仍有差距。 \\
端侧模型 & 在车端实时运行 & 算力、延迟、功耗和法规约束更强。 \\
道路验证 & 判断实际能力 & 验证成本高，长尾场景难覆盖。 \\
\bottomrule
\end{tabularx}
\end{knowledgebox}

\subsection{本章小结}

VLA 在这期里不是流行词，而是一种产品系统：它要从云端机座、世界模型、动作训练、端侧部署到道路验证全部打通。只有这样，AI 才能从“回答问题”进入“驾驶车辆”。

